\documentclass[12pt, a4paper, oneside]{article}
\usepackage{geometry}
\geometry{
    a4paper,
    top=2.5cm,
    bottom=2.5cm,
    left=2.5cm,
    right=2.5cm
}
\usepackage{amsmath}    
\usepackage{amssymb}    
\usepackage{amsfonts}   
\usepackage{amsthm}     
\usepackage{mathtools}  
\usepackage{bm}         
\usepackage{mathrsfs}

\title{\textbf{Notes}}
\author{Lin}
\date{\today}
\begin{document}
    \maketitle
    \newpage
    \newtheorem{defn}{Definition}[section]
\theoremstyle{plain}
\newtheorem{thm}[defn]{Theorem}

    $B\subseteq C \implies \mathscr{P}(B) \subseteq \mathscr{P}(C) $
    \begin{proof}

        By def $\mathscr{P}(B) = \left\{ x \mid x \subseteq B \right\}$

        $\forall x, x \in \mathscr{P}(B)$, Then, $x \subseteq B \subseteq C$

        $\implies x, x \subseteq C$
    
        $\implies x \in \mathscr{P}(C)$ By def $\mathscr{P} C = \left\{ x \mid x \subseteq C \right\}$

        Thus, $\mathscr{P}(B) \subseteq \mathscr{P}(C)$

    \end{proof}

    A is a finite set,$|A|=n \implies |\mathscr{P}(A)| = 2^n$
    \begin{proof}
        
        Consider taking $i$ elements from A to build a subset.

        Then $i \in \left\{ 0, \dots, n\right\}$, and the number of that kind of subsets is $\binom{n}{i}$.

        (if $i = 0$, the subset is $\emptyset$; if $i = n$, the subset is A itself)

        $\implies |\mathscr{P}(A)| = \sum_{i = 0}^{n}\binom{n}{i}$

        $\implies |\mathscr{P}(A)| = 2^n$

    \end{proof}

    $x, y \in B \implies \left\{ \left\{ x\right\},\left\{x, y\right\} \right\} 
    \in \mathscr{P}(\mathscr{P}(B))$
    \begin{proof}
        $x, y \in B$

        $\implies \left\{x\right\}, \left\{ x, y\right\} \in \mathscr{P}(B)$

        $\implies \left\{\left\{x\right\}, \left\{x, y\right\}\right\} \subseteq \mathscr{P}(B)$

        Thus, by def $\left\{\left\{x\right\}, \left\{x, y\right\}\right\} \in \mathscr{P}(\mathscr{P}(B))$
    \end{proof}


\end{document}

